% generated by roam.py handbook export -- не редактировать, править заметки
% 20260923123518-неравенство_коши_буняковского_шварца.org (body 56b35acc942e)

\begin{theorem}[{Неравенство Коши-Буняковского-Шварца}]\label[theorem]{note:D0601ACD-8C91-4A93-B0E4-AC0D529C6A5C}
  Пусть $X$~--- \hyperref[note:1B626AA4-35C8-4C7C-9D3D-D8C26D98F40C]{векторное пространство} над $\R$ со \hyperref[note:F2942E98-6115-4600-AA3F-E0F687D9C732]{скалярным произведением}, $\norm{x} = \sqrt{\left\langle x; x \right\rangle}$.
  Тогда
  \[
  \abs{\left\langle x; y \right\rangle} \le \norm{x} \norm{y}\quad \forall x, y \in X
  \]
  причем равенство достигается тогда и только тогда, когда $x$ и $y$ линейно зависимы.
\end{theorem}

\begin{remark}
  Раскрывая модуль, получаем двустороннюю оценку $-\norm{x} \norm{y} \le \left\langle x; y \right\rangle \le \norm{x} \norm{y}$. В оптимизации чаще нужна левая часть: линейное слагаемое $\left\langle g; y \right\rangle$ не может уменьшить значение сильнее, чем на $\norm{g} \norm{y}$.
\end{remark}

\begin{proof}
  Если $y = 0$, обе части равны нулю. Пусть $y \ne 0$. Для любого $t \in \R$ по неотрицательности скалярного квадрата
  \[
  0 \le \left\langle x - ty; x - ty \right\rangle = \norm{x}^2 - 2t \left\langle x; y \right\rangle + t^2 \norm{y}^2
  \]
  Справа~--- квадратный трехчлен по $t$ с положительным старшим коэффициентом $\norm{y}^2$, который не принимает отрицательных значений. Значит, его дискриминант неположителен:
  \[
  4 \left\langle x; y \right\rangle^2 - 4 \norm{x}^2 \norm{y}^2 \le 0
  \]
  откуда $\abs{\left\langle x; y \right\rangle} \le \norm{x} \norm{y}$.

  \proofpart{Случай равенства}
  Равенство означает, что дискриминант равен нулю, то есть у трехчлена есть корень $t_0$: $\left\langle x - t_0 y; x - t_0 y \right\rangle = 0$, а значит $x = t_0 y$.
  Обратно, при $x = t_0 y$ обе части равны $\abs{t_0} \norm{y}^2$.
\end{proof}
